\chapter{Foundations and Interfaces}

This chapter fixes three interfaces used by the surgery proof: comparison for
upper forward derivatives, the local geometry of neck surgery, and the
topology of regions covered by canonical necks and caps.

\section{Forward-Difference Comparison}

For a continuous function $f\colon[a,b]\to\mathbb R$, write
$D^+f(t)\le c$ if
\[
  \limsup_{h\downarrow0}\frac{f(t+h)-f(t)}h\le c,
  \qquad t\in[a,b).
\]

\begin{lemma}[Right-Dini mean-value lemma]
\label{lem:forward-dini-mean-value}
\dcref{MT:ch2:2.22}
\source{morgan-tian}
\notready
Let $q\colon[c,d]\to\mathbb R$ be continuous.  If $q(d)>q(c)$, then
there is $s\in[c,d)$ such that $D^+q(s)>0$.
\end{lemma}

\begin{proof}
Choose $0<\eta<(q(d)-q(c))/(d-c)$ and let $s$ minimize
$q(t)-\eta(t-c)$ on $[c,d]$.  Since
\[
 q(d)-\eta(d-c)>q(c),
\]
the point $d$ is not a minimizer, so $s<d$.  For all sufficiently small
$h>0$, minimality gives $q(s+h)-q(s)\ge\eta h$.  Hence
$D^+q(s)\ge\eta>0$.
\end{proof}

\begin{theorem}[Forward-Dini comparison principle]
\label{thm:forward-dini-comparison}
\dcref{MT:ch2:2.22}
\source{morgan-tian}
\notready
Let $f\colon[a,b]\to\mathbb R$ be continuous, let
$\psi\colon[a,b]\times\mathbb R\to\mathbb R$ be of class $C^1$, and
suppose
\[
  D^+f(t)\le \psi(t,f(t))\qquad(a\le t<b).
\]
If $G\in C^1([a,b])$ satisfies $G'(t)=\psi(t,G(t))$ and
$f(a)\le G(a)$, then $f(t)\le G(t)$ for every $t\in[a,b]$.
\end{theorem}

\begin{proof}
\uses{lem:forward-dini-mean-value}
The ranges of $f$ and $G$ lie in a compact interval $J$.  The derivative
$\partial_2\psi$ is bounded on $[a,b]\times J$, so there is $L\ge0$ such
that
\[
 |\psi(t,u)-\psi(t,v)|\le L|u-v|
 \quad(t\in[a,b],\ u,v\in J).
\]
Set $u=f-G$.  Whenever $u(t)\ge0$, subtraction of the right derivative of
$G$ from the upper forward derivative of $f$ gives
\[
 D^+u(t)\le \psi(t,f(t))-\psi(t,G(t))\le Lu(t).
\]
Consequently $v(t)=e^{-L(t-a)}u(t)$ satisfies $D^+v(t)\le0$ whenever
$v(t)\ge0$.

Suppose that $v(d)>0$ for some $d$.  Since $v(a)\le0$, continuity supplies
$c<d$ with $v(c)=0$ and $v>0$ on $(c,d]$.  By
\ref{lem:forward-dini-mean-value}, some $s\in[c,d)$ satisfies
$D^+v(s)>0$.  This contradicts $v(s)\ge0$ and the inequality above.
Thus $v\le0$, and hence $f\le G$.
\end{proof}

\begin{lemma}[Forward-Dini comparison with finitely many downward jumps]
\label{lem:forward-dini-comparison-with-downward-jumps}
\dcref{MT:ch2:2.22,MT:ch18:18.9}
\source{morgan-tian}
\uses{thm:forward-dini-comparison}
\notready
Let
\[
 a=\tau_0<\tau_1<\cdots<\tau_N=b
\]
be a finite partition.  Suppose that $f\colon[a,b]\to\mathbb R$ is
continuous on every half-open interval $[\tau_j,\tau_{j+1})$, has a finite
left lower limit at each interior partition time, and satisfies
\[
 f(\tau_j)\le\liminf_{s\uparrow\tau_j}f(s)\qquad(1\le j<N).
\]
Let $\psi\in C^1([a,b]\times\mathbb R)$ and assume
\[
 D^+f(t)\le\psi(t,f(t))
 \qquad(\tau_j\le t<\tau_{j+1})
\]
for every $j$.  If $G\in C^1([a,b])$ solves
$G'=\psi(t,G)$ and $f(a)\le G(a)$, then $f(t)\le G(t)$ on all of $[a,b]$.
\end{lemma}

\begin{proof}
On a fixed interval $[\tau_j,\tau_{j+1})$, apply
\ref{thm:forward-dini-comparison} on $[\tau_j,c]$ for every
$c<\tau_{j+1}$; the assumed right-continuity and the differential inequality
give $f(c)\le G(c)$.  Letting $c\uparrow\tau_{j+1}$ gives
\[
 \liminf_{s\uparrow\tau_{j+1}}f(s)\le G(\tau_{j+1}).
\]
At the partition time the downward-jump hypothesis therefore gives
$f(\tau_{j+1})\le G(\tau_{j+1})$.  This is the initial inequality for the
next half-open interval.  Induction over the finite partition proves the
claim, including the final endpoint.
\end{proof}

\begin{lemma}[Uniform transport of time-slice maxima]
\label{lem:uniform-time-slice-transport}
\dcref{MT:ch2:2.23}
\source{morgan-tian}
\notready
Let $X$ be a set with functions $\tau,F,V\colon X\to\mathbb R$, let
$Z\subset X$, and let $m\colon[a,b]\to\mathbb R$.  Suppose that for every
$t\in[a,b]$ there is $x\in Z$ with $\tau(x)=t$ and $F(x)=m(t)$, and that
$F(y)\le m(\tau(y))$ whenever $\tau(y)\in[a,b]$.  Suppose also that there
are $\rho>0$ and transports $\Phi_h(x)$, defined for $x\in Z$,
$|h|<\rho$, and $\tau(x)+h\in[a,b]$, such that
\[
 \tau(\Phi_h(x))=\tau(x)+h,
 \qquad
 \sup_{\substack{x\in Z\\ \tau(x)+h\in[a,b]}}
 |F(\Phi_h(x))-F(x)|\longrightarrow0
 \quad(h\to0).
\]
Assume moreover that, for every $\eta>0$, there is $h_0>0$ such that
\[
 \frac{F(x)-F(\Phi_{-h}(x))}{h}\le V(x)+\eta
\]
whenever $x\in Z$, $0<h<h_0$, and the backward transport is defined.
If $\psi\colon[a,b]\times\mathbb R\to\mathbb R$ is continuous and
\[
 V(x)\le\psi(\tau(x),F(x))\qquad(x\in Z),
\]
then $m$ is continuous and $D^+m(t)\le\psi(t,m(t))$ for every $t<b$.
\end{lemma}

\begin{proof}
For nearby $s,t\in[a,b]$, transport a maximizer $x_s\in Z$ from time $s$
to time $t$.  Since its transported value is at most $m(t)$,
\[
 m(s)-m(t)\le F(x_s)-F(\Phi_{t-s}(x_s)).
\]
Interchanging $s$ and $t$ and using the uniform value control proves the
continuity of $m$.

Fix $t<b$ and choose a maximizer $x_h\in Z$ at time $t+h$.  Backward
transport gives
\[
 \frac{m(t+h)-m(t)}h
 \le \frac{F(x_h)-F(\Phi_{-h}(x_h))}{h}.
\]
For every $\eta>0$, the right-hand side is at most
$V(x_h)+\eta$, and therefore at most
$\psi(t+h,m(t+h))+\eta$, for all sufficiently small $h>0$.
Letting $h\downarrow0$ and then $\eta\downarrow0$ proves the asserted
upper forward-derivative bound.
\end{proof}

\begin{theorem}[Time-slice maximum comparison]
\label{thm:time-slice-maximum-comparison}
\dcref{MT:ch2:2.23}
\source{morgan-tian}
\notready
Let $M$ be a smooth manifold without boundary, let $\chi$ be a smooth vector
field, and let $\mathbf t\colon M\to\mathbb R$ be smooth with
$\chi(\mathbf t)=1$.  Let $F\colon M\to\mathbb R$ be smooth and suppose that,
for every $t\in[a,b]$, $F$ attains a maximum on
$\mathbf t^{-1}(t)$.  Suppose that the set
\[
 \mathcal Z=\{x\in M:\mathbf t(x)\in[a,b]\text{ and }
 F(x)=\max_{\mathbf t^{-1}(\mathbf t(x))}F\}
\]
of maximizers in the closed slab is compact.  If
$\psi\in C^1([a,b]\times\mathbb R)$ and
\[
 \chi(F)(x)\le\psi(\mathbf t(x),F(x))
 \qquad(x\in\mathcal Z),
\]
then
\[
 F_{\max}(t)=\max_{\mathbf t^{-1}(t)}F
\]
is continuous and satisfies
$D^+F_{\max}(t)\le\psi(t,F_{\max}(t))$ for $t<b$.  Hence, if
$G'=\psi(t,G)$ and $F_{\max}(a)\le G(a)$, then
$F_{\max}(t)\le G(t)$ throughout $[a,b]$.
\end{theorem}

\begin{proof}
\uses{lem:uniform-time-slice-transport,thm:forward-dini-comparison}
Let $\varphi_h$ be the local flow of $\chi$.  Compactness of $\mathcal Z$
gives a common lifetime $\rho>0$ on a neighborhood of $\mathcal Z$, and
\[
 \mathbf t(\varphi_h(x))=\mathbf t(x)+h,
 \qquad
 \frac{d}{dh}F(\varphi_h(x))=\chi(F)(\varphi_h(x)).
\]
On a compact flow tube around $\mathcal Z$, smoothness gives both uniform
continuity of $F(\varphi_h(x))$ and the uniform expansion
\[
 \frac{F(x)-F(\varphi_{-h}(x))}{h}\longrightarrow\chi(F)(x).
\]
Thus \ref{lem:uniform-time-slice-transport}, with
$Z=\mathcal Z$, $V=\chi(F)$, and $m=F_{\max}$, proves continuity and the
forward-difference inequality.  The comparison with $G$ follows from
\ref{thm:forward-dini-comparison}.
\end{proof}

\section{Controlled-Surgery Inputs}

\subsection{Necks and singular-time extension}

\begin{definition}[Epsilon-closeness, neck, and scale]
\label{def:epsilon-neck-and-scale}
\dcref{MT:ch2:2.16,MT:ch2:2.18}
\source{morgan-tian}
\notready
Fix $0<\epsilon<1/2$ and put $k=\lfloor\epsilon^{-1}\rfloor$.  A metric
$g$ is \emph{$\epsilon$-close to $g_0$ in $C^k$} on a common domain $X$ if
\[
 \sup_{x\in X}\left(
 |g-g_0|_{g_0}^2+\sum_{\ell=1}^k
 |\nabla_{g_0}^{\ell}g|_{g_0}^2\right)<\epsilon^2.
\]

An \emph{$\epsilon$-neck} centered at a point $x$ with $R(x)>0$ in a
Riemannian three-manifold
$(M,g)$ is the image $N$ of a diffeomorphism
\[
 \varphi\colon S^2\times(-\epsilon^{-1},\epsilon^{-1})\longrightarrow N
\]
with $x\in\varphi(S^2\times\{0\})$, such that
$R(x)\varphi^*g$ is $\epsilon$-close in
$C^{\lfloor1/\epsilon\rfloor}$ to $h_0+ds^2$, where $h_0$ is the round
metric of scalar curvature $1$ on $S^2$.  The slices are the neck spheres,
the slice at $s=0$ is the central sphere, and
\[
 s_N=\operatorname{pr}_2\circ\varphi^{-1},
 \qquad r_N=R(x)^{-1/2}
\]
are the axial coordinate and the scale.  Replacing $s_N$ by $-s_N$
reverses the neck.  Its positive and negative ends are respectively
\[
 E_N^+=\{s_N>0.9\epsilon^{-1}\},
 \qquad
 E_N^-=\{s_N<-0.9\epsilon^{-1}\}.
\]
\end{definition}

\begin{definition}[Maximal extension of a generalized Ricci flow]
\label{def:maximal-ricci-flow-extension}
\dcref{MT:ch11:11.22}
\source{morgan-tian}
\notready
Here a \emph{generalized Ricci flow} consists of a manifold with a
continuous time function $\mathbf t$ which is a smooth submersion on every
ordinary or one-sided flow chart, a time vector field $\chi$ smooth on those
charts and satisfying $\chi(\mathbf t)=1$, and metrics on
$\ker(d\mathbf t)$ which evolve by the Ricci-flow equation along $\chi$.
Let $(\mathcal M,G)$ be such a
flow on a terminal interval
$[T^-,T)$, identify this part of the flow with $M\times[T^-,T)$, and set
\[
 \Omega=\{x\in M:\liminf_{t\uparrow T}R(x,t)<\infty\}.
\]
Suppose that $\Omega$ is open and that $g(t)|_\Omega$ converges smoothly
on compact subsets to a metric $g(T)$.  The \emph{maximal extension to
$T$} is
\[
 \widehat{\mathcal M}
 =\mathcal M\cup_{\Omega\times[T^-,T)}
   (\Omega\times[T^-,T]).
\]
The time function, time vector field, and horizontal metrics agree on the
overlap and hence glue to a generalized Ricci flow
$(\widehat{\mathcal M},\widehat G)$.  Its terminal time-slice is
$(\Omega,g(T))$, which need not be complete.
\end{definition}

\subsection{Local surgery data and metric}

\begin{definition}[Standard surgery-cap model]
\label{def:standard-surgery-cap-model}
\dcref{MT:ch12:12.2,MT:ch12:12.3}
\source{morgan-tian}
\uses{def:epsilon-neck-and-scale}
\notready
Fix a complete rotationally symmetric metric $g_0$ on $\mathbb R^3$ with
nonnegative sectional curvature and tip $p_0$.  Fix constants
$A_0>r_0>0$ such that $g_0$ has constant sectional curvature $1/4$ on
$B_{g_0}(p_0,r_0)$ and there is an isometry
\[
 \psi\colon(S^2\times(-\infty,4],h_0+ds^2)
 \longrightarrow(\mathbb R^3\setminus B_{g_0}(p_0,A_0),g_0).
\]
Extend the cylinder coordinate to a continuous radial function by
\[
 s_1(y)=A_0+4-d_{g_0}(p_0,y).
\]
It is smooth away from $p_0$ and has round two-sphere level sets.
\end{definition}

\begin{definition}[Local neck-surgery data]
\label{def:local-neck-surgery-data}
\dcref{MT:ch13:13.1,MT:ch13:13.2}
\source{morgan-tian}
\uses{def:epsilon-neck-and-scale,def:standard-surgery-cap-model}
\notready
Let $(N,g)$ be a $\delta$-neck centered at $x_0$, let
$Q=R_g(x_0)$, and let
\[
 \rho\colon S^2\times(-\delta^{-1},\delta^{-1})\longrightarrow N
\]
be its neck chart.  Assume $\delta^{-1}>A_0+4$.  Form $\mathcal S$ by
gluing $S^2\times(-\delta^{-1},4)$ to
$B_{g_0}(p_0,A_0+4)$ through the cylindrical isometry on
$S^2\times(0,4)$.  The two axial coordinates glue to
\[
 s\colon\mathcal S\longrightarrow(-\delta^{-1},A_0+4].
\]
Fix smooth cutoffs $\alpha\colon[1,2]\to[0,1]$ and
$\beta\colon[A_0+4-r_0,A_0+4]\to[0,1]$, equal to $1$ near their left
endpoints and $0$ near their right endpoints.  Put
\[
 \eta=1-\delta,
 \qquad
 f(s)=\begin{cases}
 0,&s\le0,\\
 C_0\delta e^{-q/s},&s>0.
 \end{cases}
\]
\end{definition}

\paragraph{Metric-convention correction for the surgery factor.}
With the explicit operator-norm convention in
\ref{def:epsilon-neck-and-scale}, the quadratic-form estimate is
$Q\rho^*g\ge(1-\delta)(h_0+ds^2)$.  Thus the factor used in the
Blueprint is $\eta=1-\delta$.  Morgan--Tian's printed
$\sqrt{1-\delta}$ is the equivalent choice under their implicit squared
quadratic-form tolerance; using it with the displayed norm would reverse
the inequality.  All subsequent surgery estimates only require
$0<\eta<1$ and $\eta\to1$.

\begin{definition}[Interpolating surgery metric]
\label{def:interpolating-surgery-metric}
\dcref{MT:ch13:13.2}
\source{morgan-tian}
\uses{def:local-neck-surgery-data}
\notready
Set $A_r=A_0+4-r_0$ and $A'=A_0+4$.  On $\mathcal S$, define
\[
\widehat g=
\begin{cases}
 e^{-2f}Q\rho^*g,&s\le1,\\
 e^{-2f}\bigl(\alpha Q\rho^*g+(1-\alpha)\eta g_0\bigr),
     &1\le s\le2,\\
 e^{-2f}\eta g_0,&2\le s\le A_r,\\
 \bigl(\beta e^{-2f}+(1-\beta)e^{-2f(A')}\bigr)\eta g_0,
     &A_r\le s\le A'.
\end{cases}
\]
The metric at the original scale is
\[
 \widetilde g=Q^{-1}\widehat g.
\]
\end{definition}

\begin{lemma}[Smoothness of the interpolating surgery metric]
\label{lem:interpolating-surgery-metric-smooth}
\dcref{MT:ch13:13.2}
\source{morgan-tian}
\notready
The piecewise tensor $\widehat g$ of
\ref{def:interpolating-surgery-metric} is a smooth Riemannian metric on
$\mathcal S$.  It equals $Q\rho^*g$ on $s\le0$.
\end{lemma}

\begin{proof}
\uses{def:interpolating-surgery-metric}
The function $e^{-q/s}$, extended by zero for $s\le0$, is smooth and flat
at $s=0$.  Near $s=1$ and $s=2$, the cutoff $\alpha$ is constant with the
value required by the adjacent formula, and the same holds for $\beta$
near $s=A_r$ and $s=A'$.  Each middle tensor is a positive linear
combination of positive-definite metrics.  Near the tip, the last formula
is the constant multiple $e^{-2f(A')}\eta g_0$, so the nonsmoothness of the
radial coordinate at $p_0$ causes no loss of smoothness.  Finally $f=0$ on
$s\le0$, which gives the asserted agreement with the normalized neck.
\end{proof}

\subsection{Distance, curvature, and convergence}

\begin{lemma}[Nonexpansion of the standard radial collapse]
\label{lem:standard-cap-radial-collapse-nonexpanding}
\dcref{MT:ch12:12.37,MT:ch13:13.1}
\source{morgan-tian}
\notready
In the data of \ref{def:local-neck-surgery-data}, define
$\xi\colon N\to\mathbb R^3$ by preserving the angular coordinate,
sending $s<A_0+4$ to the point at $g_0$-distance $A_0+4-s$ from $p_0$,
and collapsing $s\ge A_0+4$ to $p_0$.  For all sufficiently small
$\delta$,
\[
 \xi\colon(N,Qg)\longrightarrow(\mathbb R^3,\eta g_0)
\]
is distance nonincreasing.
\end{lemma}

\begin{proof}
\uses{def:local-neck-surgery-data}
The $\delta$-neck estimate and the convention-corrected
$\eta=1-\delta$ give
\[
 Q\rho^*g\ge\eta(h_0+ds^2)
\]
as quadratic forms.  Write $g_0=dr^2+u(r)^2h_0$ in radial coordinates.
The standard cap has $|dr/ds|\le1$ and $u(r)\le1$.  Thus, for
$v=a\partial_s+w$ with $w\perp\partial_s$,
\[
 |d\xi(v)|_{g_0}^2
 =a^2|dr/ds|^2+u(r)^2|w|_{h_0}^2
 \le a^2+|w|_{h_0}^2.
\]
On the collapsed region $d\xi=0$.  Hence
$|d\xi(v)|_{\eta g_0}\le|v|_{Qg}$ almost everywhere.  Integrating along
piecewise $C^1$ curves and taking infima proves the claim.
\end{proof}

\begin{lemma}[The radial surgery map is nonexpanding]
\label{lem:radial-surgery-map-nonexpanding}
\dcref{MT:ch13:13.2}
\source{morgan-tian}
\notready
Regard the map $\xi$ of
\ref{lem:standard-cap-radial-collapse-nonexpanding} as a map from $N$
to the surgery manifold $\mathcal S$.  Then
\[
 \xi\colon(N,g)\longrightarrow(\mathcal S,\widetilde g)
\]
is distance nonincreasing.
\end{lemma}

\begin{proof}
\uses{def:interpolating-surgery-metric,lem:standard-cap-radial-collapse-nonexpanding}
On $s\le0$ the normalized source and target metrics agree.  On
$0\le s\le1$, the target metric is $e^{-2f}Q\rho^*g\le Q\rho^*g$.
On $1\le s\le2$, the neck estimate gives
$\eta g_0\le Q\rho^*g$, so the convex interpolation, followed by the
factor $e^{-2f}$, is again no larger than $Q\rho^*g$.  On the remaining
cap, $\widehat g\le\eta g_0$, and
\ref{lem:standard-cap-radial-collapse-nonexpanding} applies.  Therefore
$\xi\colon(N,Qg)\to(\mathcal S,\widehat g)$ is nonexpanding.  Multiplying
both metrics by $Q^{-1}$ gives the statement.
\end{proof}

\begin{lemma}[A null-homotopic embedded sphere separates]
\label{lem:nullhomotopic-embedded-sphere-separates}
\dcref{MT:ch15:15.12}
\source{morgan-tian}
\notready
Let $M$ be a connected closed smooth $3$-manifold and let
$S\subset M$ be a smoothly embedded two-sphere.  If the inclusion
$S\hookrightarrow M$ is null-homotopic, then $M\setminus S$ is disconnected;
equivalently, $S$ separates $M$.
\end{lemma}

\begin{proof}
The sphere is two-sided, so a transverse loop has a mod-$2$ intersection number
with $S$.  If $S$ were nonseparating, choose a loop crossing from one side to
the other and close it through a product collar; its mod-$2$ intersection with
$S$ is one.  This intersection number is invariant under homotopy of the
loop and is the evaluation of the mod-$2$ Poincare-dual class of $S$.  A
null-homotopic inclusion makes that dual class zero, so every loop must have
intersection number zero, a contradiction.  Hence $S$ separates.  The same
argument works without orientability because it uses mod-$2$ coefficients.
\end{proof}

\begin{lemma}[Conformal curvature comparison on the neck side]
\label{lem:conformal-surgery-curvature-comparison}
\dcref{MT:ch13:13.5,MT:ch13:13.6,MT:ch13:13.7,MT:ch13:13.8,MT:ch13:13.9,MT:ch13:13.10,MT:ch13:13.11,MT:ch13:13.12}
\source{morgan-tian}
\notready
Let $h$ be a metric $C^{\lfloor1/\delta\rfloor}$-close to
$h_0+ds^2$ on $S^2\times I$, and put $\widehat h=e^{-2f}h$ with
$f=C_0\delta e^{-q/s}$ for $s>0$.  Choose $K$ so that every curvature
eigenvalue of a scale-one $\delta$-neck is at least $-K\delta$.
There are universal choices $q\gg(A_0+4)^2$, $C_0=2Ke^q$, and
$\delta_*>0$ such that, whenever $0<\delta\le\delta_*$ and $0<s<4$,
the scalar curvature of
$\widehat h$ is at least that of $h$, and the least eigenvalue of its
curvature operator is strictly larger than that of $h$.
\end{lemma}

\begin{proof}
In coordinates with $x^0=s$, the derivatives of $f$ satisfy
\begin{align*}
 f_0&=\frac q{s^2}f,\\
 f_{00}&=\left(\frac{q^2}{s^4}-\frac{2q}{s^3}\right)f
          +\frac q{s^2}f\,O(\delta),\\
 f_{ij}&=\frac q{s^2}f\,O(\delta)\quad(i,j)\ne(0,0).
\end{align*}
Substitution in the conformal curvature formula gives
\[
 R_{\widehat h}=e^{2f}\left(
 R_h+4\left(\frac{q^2}{s^4}-\frac{2q}{s^3}\right)f
 -2\frac{q^2}{s^4}f^2+\frac q{s^2}f\,O(\delta)\right).
\]
Choose $q$ so that $q^2/s^4$ dominates $q/s^3$ on
$(0,A_0+4]$ while $q^2s^{-4}e^{-q/s}$ remains uniformly small.  Then
choose $\delta_*$ so that $f$, the $O(\delta)$ term, and the quadratic
term in $f$ are dominated by the positive linear term.  This proves
$R_{\widehat h}\ge R_h$.

After Gram--Schmidt and diagonalization of the small curvature block, let
$\lambda\ge\mu\ge\nu$ be the eigenvalues of $h$.  Neck closeness gives
$|\lambda-1/2|+|\mu|+|\nu|\le K\delta$.  Symmetric-matrix perturbation
applied to the same conformal formula yields
\begin{align*}
 \left|\lambda'-e^{2f}\left(\lambda-\frac{q^2}{s^4}f^2\right)\right|
   &\le A\frac{q^2}{s^4}f\delta,\\
 \left|\mu'-e^{2f}\left(\mu+\left(\frac{q^2}{s^4}-\frac{2q}{s^3}\right)f\right)\right|
   &\le A\frac{q^2}{s^4}f\delta,\\
 \left|\nu'-e^{2f}\left(\nu+\left(\frac{q^2}{s^4}-\frac{2q}{s^3}\right)f\right)\right|
   &\le A\frac{q^2}{s^4}f\delta.
\end{align*}
For the choices above,
\[
 \mu'\ge e^{2f}\left(\mu+\frac{q^2}{2s^4}f\right),
 \qquad
 \nu'\ge e^{2f}\left(\nu+\frac{q^2}{2s^4}f\right).
\]
The largest eigenvalue remains $1/2+o(1)$, while both small eigenvalues
increase strictly.  Hence the least eigenvalue increases.
\end{proof}

\begin{lemma}[Curvature and pinching on the retained neck]
\label{lem:retained-neck-surgery-pinching}
\dcref{MT:ch13:13.13,MT:ch13:13.14}
\source{morgan-tian}
\notready
For the universal choices in
\ref{lem:conformal-surgery-curvature-comparison} and all sufficiently
small $\delta$, the metric $\widetilde g$ has positive sectional curvature
on $1\le s<4$.  Moreover, if the original neck satisfies, for some
$t\ge0$ and $X_g=\max(0,-\nu_g)$,
\[
 R_g\ge-\frac6{1+4t},
 \qquad
 R_g\ge2X_g\bigl(\log X_g+\log(1+t)-3\bigr)
 \quad(X_g>0),
\]
then the same inequalities hold for $\widetilde g$ on $s<4$.
\end{lemma}

\begin{proof}
\uses{def:interpolating-surgery-metric,lem:conformal-surgery-curvature-comparison}
On $1\le s<4$, the function
\[
 \frac{q^2f}{2s^4}
\]
is increasing, and at $s=1$ it equals
$q^2C_0e^{-q}\delta/2=q^2K\delta>K\delta$.  The eigenvalue estimates in
\ref{lem:conformal-surgery-curvature-comparison} therefore make both
small eigenvalues positive; the largest is greater than $1/4$ for small
$\delta$.

On $s\le0$ the metric is unchanged.  On $0<s<1$, the base metric in
\ref{def:interpolating-surgery-metric} is $Q\rho^*g$, so
\ref{lem:conformal-surgery-curvature-comparison}, followed by rescaling,
gives $R_{\widetilde g}\ge R_g$ and
$X_{\widetilde g}\le X_g$ at corresponding points.  If
$X_g\ge e^3/(1+t)$, the function
$u\mapsto2u(\log u+\log(1+t)-3)$ is increasing, so both pinching
inequalities pass to the new metric.  If $X_g<e^3/(1+t)$, the same is true
of $X_{\widetilde g}$ and the right-hand side is negative, whereas a
sufficiently round neck has nonnegative scalar curvature.  This proves the
claim on $s<1$.  On $1\le s<4$, the already established positive
sectional curvature gives $R_{\widetilde g}>0$ and
$X_{\widetilde g}=0$, so both pinching inequalities hold there as well.
\end{proof}

\begin{lemma}[Curvature and pinching on the standard-cap side]
\label{lem:standard-cap-surgery-curvature}
\dcref{MT:ch13:13.2}
\source{morgan-tian}
\notready
For the same universal choices and all sufficiently small $\delta$, the
metric $\widetilde g$ has positive sectional curvature on
$4\le s\le A_0+4$.  It satisfies the two curvature-pinching inequalities
in \ref{lem:retained-neck-surgery-pinching} there.
\end{lemma}

\begin{proof}
\uses{def:standard-surgery-cap-model,def:interpolating-surgery-metric,lem:retained-neck-surgery-pinching}
On $4\le s\le A_0+4-r_0$, the metric is a conformal multiple of
$\eta g_0$.  The standard metric has nonnegative curvature.  In its radial
warped-product coordinates, the conformal curvature formula has positive
leading terms $q^2f/s^4$ in both the radial and tangential sectional
curvatures.  Since $g_0$ is fixed, choosing $q$ after $g_0$ makes these
terms dominate all bounded lower-order coefficients; decreasing $\delta$
then absorbs the quadratic terms in $f$.  On
$A_0+4-r_0\le s\le A_0+4$, the metric converges smoothly as
$\delta\to0$ to $g_0$, which has constant sectional curvature $1/4$ on
this region.  Positivity is open in the $C^2$ topology, so it persists for
small $\delta$.  Positive sectional curvature gives $X=0$ and positive
scalar curvature, which implies both pinching inequalities.
\end{proof}

\begin{lemma}[Convergence of the surgery cap to the standard cap]
\label{lem:local-surgery-standard-cap-convergence}
\dcref{MT:ch13:13.2}
\source{morgan-tian}
\notready
For every $\delta''>0$ there is $\delta_1'(\delta'')>0$ such that, if
$\delta\le\delta_1'(\delta'')$, then $\widehat g$ on
$B_{\widehat g}(p_0,(\delta'')^{-1})$ is $\delta''$-close in
$C^{\lfloor1/\delta''\rfloor}$ to $g_0$ on the corresponding standard
ball.
\end{lemma}

\begin{proof}
\uses{def:interpolating-surgery-metric}
Fix $k=\lfloor1/\delta''\rfloor$ and the indicated standard ball.  The
cutoffs have fixed derivatives, $\eta\to1$, and every derivative through
order $k$ of $f$ tends uniformly to zero as $\delta\to0$.  On the portion
coming from the neck, the normalized metric $Q\rho^*g$ is
$\delta$-close to the product cylinder, which is exactly $g_0$ outside the
standard cap.  Thus every piece of $\widehat g$ converges to $g_0$ in
$C^k$, uniformly across the cutoff regions.  Smooth convergence also makes
the two metric balls correspond after decreasing $\delta$, proving the
stated closeness.
\end{proof}

\begin{theorem}[Local surgery on a delta-neck]
\label{thm:local-neck-surgery}
\dcref{MT:ch13:13.2}
\source{morgan-tian}
\notready
There are universal constants $C_0,q,R_0<\infty$ and $\delta_0'>0$ such
that, whenever $Q=R_g(x_0)\ge R_0$ and
$0<\delta\le\delta_0'$, the metric $\widetilde g$ of
\ref{def:interpolating-surgery-metric} has the following properties.
\begin{enumerate}
\item Every curvature-pinching inequality in
      \ref{lem:retained-neck-surgery-pinching} that holds on the original
      neck holds everywhere on $\mathcal S$.
\item The sectional curvature is positive on
      $1\le s\le A_0+4$.
\item The radial collapse map
      $\xi\colon(N,g)\to(\mathcal S,\widetilde g)$ is distance
      nonincreasing.
\item The normalized pointed cap converges smoothly on bounded balls,
      \[
        (\mathcal S,\widehat g,p_0)\longrightarrow
        (\mathbb R^3,g_0,p_0),
      \]
      as $\delta\to0$, in the quantitative sense of
      \ref{lem:local-surgery-standard-cap-convergence}.
\end{enumerate}
\end{theorem}

\begin{proof}
\uses{def:interpolating-surgery-metric,lem:interpolating-surgery-metric-smooth,lem:radial-surgery-map-nonexpanding,lem:retained-neck-surgery-pinching,lem:standard-cap-surgery-curvature,lem:local-surgery-standard-cap-convergence}
Choose the common $q$ and $C_0$ required by the curvature lemmas, and then
choose $\delta_0'$ small enough for all cited lemmas.  The two curvature
lemmas cover respectively $s<4$ and $s\ge4$, proving the first two items.
The third item is \ref{lem:radial-surgery-map-nonexpanding}, and the fourth
is \ref{lem:local-surgery-standard-cap-convergence}.  Smoothness is supplied
by \ref{lem:interpolating-surgery-metric-smooth}.
\end{proof}

\section{Canonical Neck and Cap Topology}

\subsection{Canonical regions and neck overlap}

\begin{definition}[Canonical cap]
\label{def:canonical-cap}
\dcref{MT:ch9:9.72}
\source{morgan-tian}
\uses{def:epsilon-neck-and-scale}
\notready
For $C<\infty$, a \emph{$(C,\epsilon)$-cap} is an open submanifold
$\mathcal C$ of a smooth boundaryless Riemannian three-manifold with an open
$\epsilon$-neck $N\subset\mathcal C$ whose complement in $\mathcal C$ is
compact, such that:
\begin{enumerate}
\item $\mathcal C$ is diffeomorphic to an open three-ball or to
      $\mathbb{RP}^3$ minus a point;
\item $\overline Y=\mathcal C\setminus N$ is a compact submanifold with
      boundary, its interior $Y$ is the \emph{core}, and
      $\partial\overline Y$ is a central sphere of an $\epsilon$-neck;
\item $R>0$ on $\mathcal C$,
      $\operatorname{diam}(\mathcal C)
       <C(\sup_{\mathcal C}R)^{-1/2}$.
\item $\sup_{x,y\in\mathcal C}R(y)/R(x)<C$ and
      $\operatorname{Vol}(\mathcal C)
      <C(\sup_{\mathcal C}R)^{-3/2}$.
\item For $y\in Y$, let $r_y>0$ satisfy
      $\sup_{B(y,r_y)}R=r_y^{-2}$.  Then
      $B(y,r_y)\Subset\mathcal C$ and
      \[
       \inf_{y\in Y}\frac{\operatorname{Vol}B(y,r_y)}{r_y^3}>C^{-1}.
      \]
\item The scale-invariant derivative bounds
      \[
       \sup_{\mathcal C}\frac{|\nabla R|}{R^{3/2}}<C,
       \qquad
       \sup_{\mathcal C}
       \frac{|\Delta R+2|\operatorname{Ric}|^2|}{R^2}<C
      \]
      hold.
\end{enumerate}
The neck orientation is fixed so that the closure of its negative end meets
the core; reversing this orientation is not a new cap structure.
\end{definition}

\begin{definition}[Embedded horn frontier and cross-sectional sphere]
\label{def:embedded-horn-frontier}
\dcref{MT:ch9:9.72,MT:app:A.21}
\source{morgan-tian}
\uses{def:epsilon-neck-and-scale,def:canonical-cap}
\notready
For an embedded neck chart $\varphi_N$ as in
\ref{def:epsilon-neck-and-scale}, write
\[
 S_N(\sigma)=\varphi_N(S^2\times\{\sigma\}),\qquad
 H_N^+=\varphi_N(S^2\times(0,\epsilon^{-1})),\qquad
 H_N^-=\varphi_N(S^2\times(-\epsilon^{-1},0)).
\]
The \emph{horn frontier} of a cap is the embedded central sphere
$S_N(0)=\partial\overline Y$ of its neck end, with the side containing the
compact core designated the core side.  For any open region $U\subset M$,
\(\operatorname{Fr}_M(U)=\overline U^{\,M}\setminus U\) denotes its ambient
frontier.  Thus every frontier point below is a point of this specified
embedded sphere (or of an exposed end frontier), and every
cross-sectional sphere is one of the displayed $S_N(\sigma)$.
\end{definition}

\begin{definition}[Embedded epsilon-horn and end convention]
\label{def:epsilon-horn-terminology}
\dcref{KL:II:Definition2}
\source{kleiner-lott}
\uses{def:epsilon-neck-and-scale}
\notready
Let $(M^3,g)$ be a smooth boundaryless Riemannian three-manifold and let
$I\subset\mathbb R$ be a nonempty open interval.  An \emph{embedded
$\epsilon$-horn} is an open embedded submanifold
$\mathcal H\cong S^2\times I$ such that every point of $\mathcal H$ lies in
an $\epsilon$-neck and, for one end of $I$, the scalar curvature remains
bounded along that end while it tends to $+\infty$ along the other end.  The
chosen product coordinate is proper toward the divergent end; its compact
cross-sectional spheres are the embedded frontier spheres used for truncation.
Reversing the product coordinate exchanges the two ends but does not change
the horn.  A capped $\epsilon$-horn is obtained by adjoining one canonical cap
along the bounded end; a double $\epsilon$-horn has divergent ends at both
sides.  In a Ricci-flow terminal slice, the bounded-end condition means
bounded curvature on the terminal domain, and the divergent-end condition means
$\bar R\to+\infty$ along the end.
\end{definition}

\begin{definition}[Tubes, capped tubes, and neck fibrations]
\label{def:canonical-neck-cap-regions}
\dcref{MT:ch9:9.81,MT:app:A.21,MT:app:A.25}
\source{morgan-tian}
\uses{def:epsilon-neck-and-scale,def:canonical-cap}
\notready
An \emph{$\epsilon$-tube} is a submanifold diffeomorphic to
$S^2\times I$, for a nondegenerate interval $I$, covered by
$\epsilon$-necks whose central spheres are isotopic to the $S^2$ factors.
An open tube has no boundary.  A \emph{$C$-capped $\epsilon$-tube} is the
connected union of an open $\epsilon$-tube and one $(C,\epsilon)$-cap,
glued along an end of each.  A \emph{doubly $C$-capped
$\epsilon$-tube} has disjoint cap cores attached at both ends.  An
\emph{$\epsilon$-fibration} is a closed connected $S^2$-bundle over
$S^1$, covered by $\epsilon$-necks whose central spheres are isotopic to
the fibers.
\end{definition}

\begin{definition}[Balanced chain of epsilon-necks]
\label{def:balanced-neck-chain}
\dcref{MT:app:A.12,MT:app:A.17}
\source{morgan-tian}
\uses{def:epsilon-neck-and-scale}
\notready
A finite \emph{chain of $\epsilon$-necks} is a sequence
$N_a,\ldots,N_b$ with oriented axial coordinates such that, for each
$a\le i<b$,
\[
 \{s_{N_i}>\tfrac12\epsilon^{-1}\}\subset N_i\cap N_{i+1},
 \qquad
 \{s_{N_{i+1}}<-\tfrac12\epsilon^{-1}\}\subset N_i\cap N_{i+1},
\]
while
\[
 N_i\cap N_{i+1}\subset
 \{s_{N_i}>-\tfrac12\epsilon^{-1}\}
 \cap\{s_{N_{i+1}}<\tfrac12\epsilon^{-1}\},
\]
and $N_i\cap E_{N_a}^-=\varnothing$ for
$i>a$.  An infinite chain indexed by an interval in $\mathbb Z$ is one
whose restriction to every finite subinterval is a chain.  The chain is
\emph{balanced} if consecutive centers $x_j,x_{j+1}$ satisfy
\[
 0.99\,r_{N_j}\epsilon^{-1}
 \le d(x_j,x_{j+1})
 \le1.01\,r_{N_j}\epsilon^{-1}.
\]
\end{definition}

\begin{lemma}[Metric control in a small epsilon-neck]
\label{lem:epsilon-neck-metric-control}
\dcref{MT:app:A.2,MT:app:A.5}
\source{morgan-tian}
\notready
Write $d_N$ for the intrinsic distance induced by the restriction of the
ambient metric to $N$.  The estimates below are intrinsic to the neck; no
ambient no-shortcut assertion is being made.
For every $0<\alpha<1/8$, all sufficiently small $\epsilon$ have the
following property.  If $N$ is an $\epsilon$-neck of scale $r_N$ centered
at $x$, then scalar curvatures on $N$ lie between
$(1-\alpha)r_N^{-2}$ and $(1+\alpha)r_N^{-2}$.  Neck spheres have
diameter at most $(1+\alpha)\sqrt2\pi r_N$, and points with sufficiently
large axial separation, namely
$|s_N(q)-s_N(p)|\ge\epsilon^{-1}/100$ or
$d_N(p,q)\ge r_N\epsilon^{-1}/100$, satisfy
\[
 (1-\alpha)r_N|s_N(q)-s_N(p)|
 \le d_N(p,q)
 \le(1+\alpha)r_N|s_N(q)-s_N(p)|.
\]
In particular, every path contained in $N$ and joining its two end regions
has length at least
$\frac95(1-\alpha)r_N\epsilon^{-1}$.
\end{lemma}

\begin{proof}
\uses{def:epsilon-neck-and-scale}
After multiplying the metric by $r_N^{-2}$, the neck is arbitrarily close
in $C^2$ to the round product cylinder.  Curvature continuity gives the
scalar-curvature bounds and the diameter bound for the slices.  In the
product cylinder, a minimizing geodesic has an axial projection of speed at
most one, while its spherical projection has uniformly bounded diameter.
For an axial displacement tending to infinity with $\epsilon^{-1}$, the
spherical error is absorbed into the factor $1\pm\alpha$.  The same
argument applies to intrinsic distance in $N$ when a large intrinsic
distance, rather than a large axial displacement, is assumed, since the
spherical diameter is bounded.
Rescaling by $r_N$ proves the distance estimate.  A path joining the two end
regions has axial displacement at least $1.8\epsilon^{-1}$, which gives the
last assertion.
\end{proof}

\begin{proposition}[Overlap control for small epsilon-necks]
\label{prop:overlapping-epsilon-necks}
\dcref{MT:app:A.11}
\source{morgan-tian}
\notready
For every $0<\alpha\le10^{-2}$, all sufficiently small $\epsilon$ have
the following property.  If two $\epsilon$-necks $N,N'$ intersect, then
their scales differ by a factor in $(1-\alpha,1+\alpha)$, their neck-sphere
tangent planes make angle less than $\alpha$ on the overlap, and their
oriented axial directions agree up to sign and an error less than
$\alpha$.  A neck sphere of one neck contained in the other is isotopic
there to a neck sphere.  If the overlap contains a point in the middle
nine-tenths of one neck, then there is a point in the overlap lying in the
middle $0.96\epsilon^{-1}$ of both neck coordinates; the corresponding
cross-sectional spheres are contained in the opposite necks, isotopic in the
overlap, and the relevant overlap component is diffeomorphic to
$S^2\times(0,1)$.
\end{proposition}

\begin{proof}
\uses{def:epsilon-neck-and-scale,lem:epsilon-neck-metric-control}
At an intersection point, both normalized metrics are close to the same
round-cylinder metric.  Comparing scalar curvature there gives the scale
ratio.  On a sufficiently round cylinder the curvature operator has a
simple eigenvalue near $1/2$ whose eigenspace is the tangent plane of the
sphere factor; continuity of this eigenspace proves the tangent-plane
estimate.  Its orthogonal line gives the axial directions, determined up
to sign.

If a sphere $S'$ from $N'$ lies in $N$, projection along the axial lines of
$N$ restricts to a local diffeomorphism $S'\to S^2$.  It is a covering
because $S'$ is compact, and hence a diffeomorphism because $S^2$ is simply
connected.  Thus $S'$ is a graph over a neck sphere and is isotopic to it
along axial lines.  In a middle overlap, the metric and axial-coordinate
estimates provide such graph spheres from both neck structures.  The
isotopy between them, together with their collars, supplies the asserted
product structure on the overlap.
\end{proof}

\begin{lemma}[Separation type is constant on a connected neck-center set]
\label{lem:neck-separation-type-locally-constant}
\dcref{MT:app:A.11}
\source{morgan-tian}
\uses{def:epsilon-neck-and-scale,prop:overlapping-epsilon-necks}
\notready
For sufficiently small $\epsilon$, let $X$ be a connected subset of a
smooth boundaryless Riemannian three-manifold $M$ such that every point of
$X$ is the center of an embedded $\epsilon$-neck.  Then either every central
neck sphere arising from points of $X$ separates $M$, or none does.  More
precisely, if two centers are sufficiently close that their necks meet in
the middle nine-tenths, their central spheres are ambiently isotopic and
therefore have the same separation predicate; these neighborhoods form an
open cover of $X$.
\end{lemma}

\begin{proof}
\uses{prop:overlapping-epsilon-necks}
The middle-overlap clause of
\ref{prop:overlapping-epsilon-necks} supplies an isotopy between the two
central spheres inside the overlap.  Ambient isotopy preserves whether the
complement has one or two components.  Neck charts vary continuously with
their centers, so a sufficiently small neighborhood of each center has this
middle-overlap property; hence the separating and nonseparating loci in $X$
are both open.  Since $X$ is connected, at most one is nonempty.
\end{proof}

\subsection{Chains of necks}

\begin{lemma}[A neck chain is a product]
\label{lem:neck-chain-product}
\dcref{MT:app:A.13}
\source{morgan-tian}
\notready
Let $N_j$ be a balanced, non-returning chain of sufficiently small
$\epsilon$-necks, with $N_i\cap N_j=\varnothing$ whenever
$|i-j|\ge3$.  For an infinite chain assume the positive-end frontier
condition stated below; for a finite chain assume
the two exposed ends are distinct.  Then the union
$U=\bigcup_jN_j$ is diffeomorphic to $S^2\times(0,1)$, compatibly with the
isotopy classes of all neck spheres.
\end{lemma}

\begin{proof}
\uses{def:balanced-neck-chain,prop:overlapping-epsilon-necks}
For a finite chain, induct on the number of necks.  When the last neck is
adjoined, \ref{prop:overlapping-epsilon-necks} identifies the overlap with
$S^2\times(0,1)$ and identifies a cross-sectional sphere with the factor
sphere in both pieces.  Gluing the corresponding collars therefore extends
the product coordinate across the new neck.  The construction can be made
compatible under inclusion.  For an infinite chain, the
nonconsecutive-disjointness and the stated positive-end frontier condition
exclude returns and
accumulation at the positive end.  The compatible finite product coordinates
therefore exhaust a proper interval coordinate on $U$, giving
$U\cong S^2\times(0,1)$.
\end{proof}

\begin{lemma}[An infinite neck-chain end has no frontier]
\label{lem:infinite-neck-chain-frontier}
\dcref{MT:app:A.15}
\source{morgan-tian}
\notready
Let $N_0,N_1,\ldots$ be a chain of sufficiently small
$\epsilon$-necks in $M$, and put $U=\bigcup_{j\ge0}N_j$.  Then the
frontier of $U$ in $M$ is exactly the frontier of the negative end of
$N_0$.
\end{lemma}

\begin{proof}
\uses{def:balanced-neck-chain,lem:epsilon-neck-metric-control}
The frontier of a finite subchain consists of the two exposed end
frontiers.  Hence a limit point of $U$ represented by points in one fixed
finite subchain can lie in the frontier of $U$ only at the negative end of
$N_0$: the positive end of every finite subchain lies in the next neck.

It remains to exclude a convergent sequence $z_k\in N_{j_k}$ with
$j_k\to\infty$.  If $z_k\to z$, continuity of scalar curvature bounds
$R(z_k)$ and gives a positive lower bound for the scales of all sufficiently
late necks.  By \ref{lem:epsilon-neck-metric-control}, those necks have a
uniformly positive width.  Any path from $z_k$ to $z$ must cross either
$N_0$ or a complete sufficiently late neck, so $d(z_k,z)$ has a fixed
positive lower bound, a contradiction.  Thus the infinite end contributes
no frontier.
\end{proof}

\begin{lemma}[Extension of a balanced neck chain]
\label{lem:balanced-neck-chain-extension}
\dcref{MT:app:A.17,MT:app:A.18}
\source{morgan-tian}
\notready
For all sufficiently small $\epsilon$, let
$N_a,\ldots,N_b$ be a balanced chain whose central spheres separate $M$.
If a point in the frontier of its positive end is the center of an
$\epsilon$-neck $N$, then, after reversing $N$ if necessary,
$N_a,\ldots,N_b,N$ is a balanced chain.  The analogous assertion holds at
the negative end.
\end{lemma}

\begin{proof}
\uses{def:balanced-neck-chain,lem:epsilon-neck-metric-control,prop:overlapping-epsilon-necks}
The two cases are symmetric.  At the positive end, the new center lies in
the closure of $N_b$.  The distance estimate in
\ref{lem:epsilon-neck-metric-control} gives the balanced inequalities for
the two centers, and \ref{prop:overlapping-epsilon-necks} allows the axial
coordinate of $N$ to be oriented consistently with $N_b$.  The separating
central sphere of $N_a$ prevents $N$ from returning to the negative end of
the chain: otherwise a path through the overlaps would cross that sphere
exactly once.  Thus all defining conditions of a balanced chain remain
valid.
\end{proof}

\begin{lemma}[Finite neck-chain frontier extension dichotomy]
\label{lem:finite-neck-chain-frontier-dichotomy}
\dcref{MT:app:A.17,MT:app:A.18,MT:app:A.21}
\source{morgan-tian}
\uses{def:balanced-neck-chain,lem:epsilon-neck-metric-control,
  prop:overlapping-epsilon-necks,lem:balanced-neck-chain-extension}
\notready
Let $X$ be connected in a boundaryless smooth Riemannian three-manifold $M$,
and let $U$ be the compact closure of the union of a finite balanced chain of
sufficiently small $\epsilon$-necks.  Assume
$X\cap\operatorname{int}U\ne\varnothing$ and
$X\cap(M\setminus\overline U)\ne\varnothing$, and that every point of
$X\setminus\operatorname{int}U$ is the center of an $\epsilon$-neck whose
central sphere separates $M$.  Then $X$ meets the frontier of one exposed end
of $U$ at a neck center, and that neck extends the balanced chain after an
axial reorientation.
\end{lemma}

\begin{proof}
The sets $X\cap\operatorname{int}U$ and
$X\cap(M\setminus\overline U)$ are disjoint relatively open subsets of $X$.
Since $X$ is connected and meets both, the remaining set
$X\cap\operatorname{Fr}_M(U)$ is nonempty.  A frontier point cannot lie in
the interior of a chain neck, so it lies in the exposed end region of the
first or last neck.  The overlap and metric-control estimates in
\ref{prop:overlapping-epsilon-necks} and
\ref{lem:epsilon-neck-metric-control} identify its neck with that end and
give the scale and axial alignment hypotheses of
\ref{lem:balanced-neck-chain-extension}.  The separating-sphere assumption
is precisely the nonreturn condition in that extension lemma, so adjoining
the neck extends $U$ at the relevant end.
\end{proof}

\begin{theorem}[Separating neck centers lie in an epsilon-tube]
\label{thm:separating-neck-centers-form-tube}
\dcref{MT:app:A.12,MT:app:A.15,MT:app:A.17,MT:app:A.18,MT:app:A.19}
\source{morgan-tian}
\notready
There is $\epsilon_0>0$ such that the following holds for
$0<\epsilon\le\epsilon_0$.  Let $X$ be a connected subset of a smooth
Riemannian three-manifold $M$ without boundary.  Suppose every $x\in X$ is
the center of an embedded $\epsilon$-neck and every corresponding central
sphere separates $M$.
Then a subcollection of these necks, after reversing axial coordinates when
necessary, is a balanced chain whose union contains $X$.  Consequently
$X$ lies in an $\epsilon$-tube.
\end{theorem}

\begin{proof}
\uses{def:canonical-neck-cap-regions,lem:neck-chain-product,
  lem:infinite-neck-chain-frontier,lem:balanced-neck-chain-extension,
  lem:finite-neck-chain-frontier-dichotomy}
Start with one neck centered in $X$.  The finite-chain frontier lemma and its
balanced extension clause are the source's no-branching construction: at each
exposed end, every center of $X$ is forced into the same end overlap, and the
separating central sphere prevents a return through the opposite side.  Thus
the successive necks form one ordered chain rather than a branching tree.
Iterating the extension gives either a finite chain containing $X$, a
one-sided infinite chain, or a two-sided chain.  The infinite-chain frontier
lemma excludes a frontier at an infinite end; if an end remains finite,
connectedness and the explicit exterior hypothesis in
\ref{lem:finite-neck-chain-frontier-dichotomy} force the next extension.
  Consequently the chain union contains $X$ (or is the ambient component in the
  two-sided case).  The product lemma identifies it with
  $S^2\times(0,1)$, hence an $\epsilon$-tube.  The separating hypothesis is
  used at every extension step, while the nonseparating case is handled by the
  separate return/fibration dichotomy below.
\end{proof}

\begin{lemma}[Nonseparating neck chains: tube or closed fibration]
\label{lem:nonseparating-neck-chain-or-bundle}
\dcref{MT:app:A.20,MT:app:A.24}
\source{morgan-tian}
\uses{def:balanced-neck-chain,prop:overlapping-epsilon-necks,
  lem:epsilon-neck-metric-control,lem:neck-chain-product,
  lem:infinite-neck-chain-frontier}
\notready
For sufficiently small $\epsilon$, let $X$ be a connected subset of a
boundaryless Riemannian three-manifold $M$ such that every point of $X$ is the
center of an embedded $\epsilon$-neck and no corresponding central sphere
separates $M$.  Then either
\begin{enumerate}
\item there is a balanced chain of these necks, with axial coordinates
      possibly reversed, whose union $U\cong S^2\times(0,1)$ contains $X$; or
\item $X$ is contained in a connected component $M_0\subset M$ all of whose
      points are $\epsilon$-neck centers, and $M_0$ is a closed
      $S^2$-bundle over $S^1$ (an $\epsilon$-fibration).
\end{enumerate}
\end{lemma}

\begin{proof}
Start with one neck centered in $X$ and extend a finite balanced chain whenever
an exposed end meets another center of $X$.  The overlap proposition and the
metric-control estimate give the next balanced neck unless it meets the
negative end of the first neck; that is the only new possibility because the
central spheres are nonseparating.  If no return occurs, the extension process
is finite or one- or two-sided infinite.  The infinite-chain frontier lemma,
together with connectedness of $X$, then shows that the resulting chain union
contains $X$, and the neck-chain product lemma gives
$U\cong S^2\times(0,1)$.

Suppose instead that the chain returns to its first negative end.  If the
return overlap reaches the middle nine-tenths of that first neck, the overlap
proposition makes the finite closed chain a whole connected component.  If the
return lies only in the far negative end, use connectedness of $X$ exactly as
in the source's Case~II: either replace the first center by a center in the
positive part of the first neck, producing a chain containing $X$, or retain a
center in the far negative part, in which case the closed chain is a whole
component.  In the latter case its universal cover is the bi-infinite neck
chain $S^2\times\mathbb R$; the deck action preserves the neck foliation and
is a nontrivial translation on the interval coordinate, so the quotient is a
closed $S^2$-bundle over $S^1$.  This is the return/fibration argument in
Appendix~A.20; the cap-gluing return claim used in the far-negative case is
Appendix~A.24, while the containment alternatives are recorded separately in
A.21.

The printed statement of Appendix~A.19 has the same nonseparating hypothesis
but its proof invokes the separating-sphere extension lemma A.17--A.18; that
polarity mismatch is why this is a repaired, non-literal import.  The closed
product $S^2\times S^1$ is already a counterexample to the printed
unconditional tube conclusion.  The Blueprint therefore uses the A.20
return/fibration classification (and only its A.24 cap-gluing subclaim where
appropriate) and retains the fibration alternative.
\end{proof}

\begin{lemma}[All-neck components are tubes or sphere bundles]
\label{lem:all-neck-component-tube-or-bundle}
\dcref{MT:app:A.20}
\source{morgan-tian}
\uses{thm:separating-neck-centers-form-tube,lem:neck-chain-product}
\notready
For sufficiently small $\epsilon$, let $Y$ be a connected Riemannian
three-manifold without boundary such that every point of $Y$ is the center of
an embedded $\epsilon$-neck.  Then either $Y$ is diffeomorphic to
$S^2\times(0,1)$ and is an $\epsilon$-tube, or $Y$ is diffeomorphic to an
$S^2$-bundle over $S^1$.
\end{lemma}

\begin{proof}
Lift the neck cover to the universal cover $\widetilde Y$.  Every lifted neck
center is still a neck center.  A lifted central sphere separates
$\widetilde Y$: otherwise a path crossing it once, closed using a collar,
would give a nontrivial loop in the simply connected cover.  The separating
neck theorem and the neck-chain product lemma therefore identify
$\widetilde Y$ with $S^2\times(0,1)$, with its neck spheres as the interval
fibers.  If the deck group is trivial, $Y$ is the open tube.  Otherwise
proper discontinuity and the disjointness of sufficiently small translated
neck charts force the induced action on the interval coordinate to be a
nontrivial cyclic translation action; its quotient is a closed
$S^2$-bundle over $S^1$.  This is the source's Appendix~A.20 universal-cover
dichotomy and does not use the polarity-ambiguous printed A.19 proof.
\end{proof}

\subsection{Caps and classification}

\begin{lemma}[Uniform scale in a nested cap chain]
\label{lem:nested-cap-scale-comparison}
\dcref{MT:app:A.22}
\source{morgan-tian}
\uses{def:canonical-cap,prop:overlapping-epsilon-necks,
  lem:epsilon-neck-metric-control}
\notready
Fix $C$ and sufficiently small $\epsilon$.  Let
$C_0\subset C_1\subset\cdots$ be canonical caps whose successive core
closures meet the preceding frontier.  If $Q_0$ is the scalar
curvature at a fixed point of the first core, then every neck scale exposed by
the chain is at least $c(C,\epsilon)Q_0^{-1/2}$, and every noncontradictory
enlargement increases the distance from that point to the complement by at
least $c(C,\epsilon)Q_0^{-1/2}$.
\end{lemma}

\begin{proof}
The cap curvature-ratio bound compares each core scale with the scale at the
point where the preceding cap meets its frontier.  The overlap proposition
then compares that scale with the next neck scale, while neck metric control
converts it to an axial-width lower bound.  Induction along the nested chain
gives the displayed uniform bound and distance increment.  This is the
quantitative estimate in the proof of Morgan--Tian Appendix~A.22; it is
recorded separately so the no-infinite-chain conclusion does not silently
assume a cross-cap curvature bound.
\end{proof}

\begin{lemma}[No infinite nested chain of canonical caps]
\label{lem:no-infinite-nested-canonical-caps}
\dcref{MT:app:A.22}
\source{morgan-tian}
\uses{def:canonical-cap,lem:nested-cap-scale-comparison}
\notready
For fixed $C$ and sufficiently small $\epsilon$, there is no chain of
$(C,\epsilon)$-caps
\[
 C_0\subset C_1\subset\cdots
\]
such that the closure of the core of $C_i$ meets
$\operatorname{Fr}_M(C_{i-1})$ for every $i\ge1$.
\end{lemma}

\begin{proof}
\uses{def:canonical-cap,lem:epsilon-neck-metric-control,prop:overlapping-epsilon-necks,
  lem:nested-cap-scale-comparison}
Fix $x_0\in C_0$ and put $Q_0=R(x_0)$.  Let $S_i'$ be the boundary of
the core of $C_i$, realized as the central sphere of a neck $N_i'$, and
let $N_i$ be the neck end of $C_i$.  The uniform scale contract
\ref{lem:nested-cap-scale-comparison} gives every $N_i$ scale at least
$c(C,\epsilon)Q_0^{-1/2}$.

If $S_i'$ is disjoint from $C_{i-1}$, the component of
$M\setminus S_i'$ containing $C_{i-1}$ must be the core side; otherwise
$C_{i-1}\cap C_i$ would lie entirely in the neck end, contrary to
$C_{i-1}\subset C_i$.  Hence the core of $C_i$ contains $C_{i-1}$, and
\ref{lem:nested-cap-scale-comparison} gives
\[
 d(x_0,M\setminus C_i)-d(x_0,M\setminus C_{i-1})
  \ge c(C,\epsilon)Q_0^{-1/2}
\]
for the same $c(C,\epsilon)$.  If $S_i'\subset C_{i-1}$, the component of its
complement whose closure lies in $C_{i-1}$ cannot be the core side, since
that core meets $\operatorname{Fr}(C_{i-1})$.  It would therefore contain
the neck end of $C_i$ while having only one boundary sphere, contradicting
the two frontier components of a neck end.

In the remaining case $S_i'$ crosses $N_{i-1}$.
By \ref{prop:overlapping-epsilon-necks}, the cylinder directions almost
agree or almost oppose.  Separation forces the agreeing orientation, so
the central sphere of $N_i$ lies beyond the positive end of $N_{i-1}$.
  The same comparison lemma gives the displayed increase.  Thus every
noncontradictory step raises the distance to the complement by a fixed
positive amount.  But \ref{def:canonical-cap} bounds
$\operatorname{diam}(C_i)$ by $CQ_0^{-1/2}$.  This is impossible for
arbitrarily large $i$.
\end{proof}

\begin{lemma}[An isotopic cap-gluing sphere determines a component]
\label{lem:isotopic-cap-gluing-sphere}
\dcref{MT:app:A.24}
\source{morgan-tian}
\notready
Let $X\subset M$ be connected.  Let $\widetilde C$ be a cap or a finite
capped tube containing a point of $X$, and let $C'$ be a
$(C,\epsilon)$-cap whose core contains a point of
$X\cap\operatorname{Fr}_M(\widetilde C)$.  Let $N'$ be the neck end of
$C'$, oriented away from its core.  If a two-sphere
$\Sigma\subset N'\cap\widetilde C$ lies in the closed positive half of
$N'$ and is isotopic in $N'$ to its central sphere, then
$\widetilde C\cup C'$ is a component of $M$ containing $X$.
\end{lemma}

\begin{proof}
\uses{def:canonical-cap,def:canonical-neck-cap-regions}
The sphere $\Sigma$ separates $\widetilde C$ into a relatively compact
side $A$ and an end side.  It also separates $C'$; let $A'$ be the side
with compact closure, which contains the core of $C'$ because $\Sigma$ is
isotopic to the central sphere of $N'$.  Both $A$ and $A'$ have frontier
$\Sigma$.  They cannot lie on the same side of $\Sigma$, since then the
core of $C'$ would lie in $\widetilde C$, contrary to its meeting the
frontier of $\widetilde C$.  Thus they lie on opposite sides, and their
closures form an ambient component.  This component is
$\widetilde C\cup C'$.  It meets $X$, so connectedness of $X$ places all
of $X$ in it.
\end{proof}

\begin{lemma}[Adjacent canonical caps form an ambient component]
\label{lem:adjacent-canonical-caps-form-component}
\dcref{MT:app:A.23}
\source{morgan-tian}
\notready
Let $X\subset M$ be connected, and let $\widetilde C$ be obtained from a
$(C,\epsilon)$-cap $C_0$ by adjoining a finite balanced chain to its neck
end.  Suppose $\widetilde C$ contains a point of $X$.  Let $C'$ be a
$(C,\epsilon)$-cap whose core contains a point of
$X\cap\operatorname{Fr}_M(\widetilde C)$, and assume
$C_0\nsubseteq C'$.  Then $\widetilde C\cup C'$ is a component of $M$
containing $X$.
\end{lemma}

\begin{proof}
\uses{def:canonical-cap,def:balanced-neck-chain,prop:overlapping-epsilon-necks,lem:isotopic-cap-gluing-sphere}
Let $S'$ be the boundary of the core of $C'$ and $N'$ its neck end.
If $S'\subset\widetilde C$, apply
\ref{lem:isotopic-cap-gluing-sphere} with $\Sigma=S'$.  If $S'$ is
disjoint from $\widetilde C$, the side of $S'$ containing
$\widetilde C$ cannot be the core side, since that would imply
$C_0\subset C'$.  Hence the neck $N'$ meets the last neck of the balanced
chain.  By \ref{prop:overlapping-epsilon-necks}, their overlap contains a
sphere $\Sigma$ isotopic to $S'$ and lying in the positive half of $N'$.
Apply \ref{lem:isotopic-cap-gluing-sphere}.

Finally suppose $S'$ crosses the last neck of $\widetilde C$.  The scale,
diameter, and orientation conclusions of
\ref{prop:overlapping-epsilon-necks} again give a common sphere
$\Sigma$ in the positive half of $N'$.  If the cap core and
$\widetilde C$ lie on opposite sides of $\Sigma$, apply
\ref{lem:isotopic-cap-gluing-sphere}.  If they lie on the same side, the
core of $C'$ contains $\widetilde C$ up to its last half-neck; the overlap
then supplies that half-neck as well, forcing $C_0\subset C'$, a
contradiction.  These cases are exhaustive.
\end{proof}

\begin{lemma}[Frontier localization for a capped finite neck chain]
\label{lem:capped-chain-exposed-frontier}
\dcref{MT:app:A.17,MT:app:A.21}
\source{morgan-tian}
\uses{def:canonical-cap,def:balanced-neck-chain,
  lem:neck-chain-product,prop:overlapping-epsilon-necks,
  lem:balanced-neck-chain-extension}
\notready
Let $C_0$ be a $(C,\epsilon)$-cap and let
$N_a,\ldots,N_b$ be a finite balanced chain whose negative end is attached to
$C_0$, with the positive end of $N_b$ the only uncapped end.  More precisely,
after the cap orientation of \ref{def:canonical-cap} is fixed, assume that
$E_{N_a}^-\subset C_0$ and that all other cap/neck intersections are the
overlaps in \ref{def:balanced-neck-chain}.  Assume that
\[
U=\overline{C_0\cup N_a\cup\cdots\cup N_b}
\]
is compact.  Then every point of
$\operatorname{Fr}_M(U)$ lies in the exposed positive-end region of $N_b$.
Consequently, if such a frontier point is the center of an
$\epsilon$-neck, that neck can be axially oriented to extend the chain.
\end{lemma}

\begin{proof}
The neck-chain product lemma identifies the union of the overlapping necks
with $S^2\times I$ and shows that its only frontier components are the two
axial end fibers.  The negative end fiber is contained in the cap core by the
attachment hypothesis, so it is not in the frontier of the capped union.
Thus every ambient frontier point lies at the positive end of $N_b$.  The
overlap proposition compares any neck centered there with $N_b$, and the
balanced-chain extension lemma supplies the orientation and the new chain
conditions.  This is the exposed-end localization used in Appendix~A.21.
\end{proof}

\begin{lemma}[Finite cap-chain frontier dichotomy]
\label{lem:finite-cap-chain-frontier-dichotomy}
\dcref{MT:app:A.21}
\source{morgan-tian}
\uses{def:canonical-neck-cap-regions,def:balanced-neck-chain,
  prop:overlapping-epsilon-necks,lem:capped-chain-exposed-frontier}
\notready
Let $X$ be connected and let $U$ be the compact closure of a connected union
of one $(C,\epsilon)$-cap and a finite balanced chain of $\epsilon$-necks,
attached at the negative end with only the positive end exposed as in
\ref{lem:capped-chain-exposed-frontier}.
Assume $X\cap\operatorname{int}U\ne\varnothing$ and
$X\cap(M\setminus\overline U)\ne\varnothing$, and every point of
$X\setminus\operatorname{int}U$ is either an $\epsilon$-neck center or lies in
the core of a $(C,\epsilon)$-cap.  If $X\not\subset U$, then
$X\cap\operatorname{Fr}_M(U)$ contains either a neck center whose neck extends
the balanced chain or a point in the core of a second cap.
\end{lemma}

\begin{proof}
If $X$ meets both $\operatorname{int}U$ and $M\setminus U$, then a connected
set cannot avoid $\operatorname{Fr}_M(U)$: otherwise its intersections with
$\operatorname{int}U$ and with $M\setminus\overline U$ would be a separation
of $X$.  Thus $X\cap\operatorname{Fr}_M(U)\ne\varnothing$.  By the
covering hypothesis it is either a neck center or belongs to a cap core.  In
the neck case, \ref{lem:capped-chain-exposed-frontier} localizes the frontier
point to the positive end of the last neck and supplies the overlap and
balanced extension.  In the cap case it is, by definition, a distinct second
cap core meeting the frontier.  The
two cases are exhaustive.  This is the finite-chain frontier step in
Morgan--Tian Appendix A.21, stated without assuming the classification
conclusion.
\end{proof}

\begin{lemma}[Maximal cap selection along a connected covered set]
\label{lem:maximal-cap-selection}
\dcref{MT:app:A.21,MT:app:A.22,MT:app:A.23}
\source{morgan-tian}
\uses{def:canonical-cap,lem:no-infinite-nested-canonical-caps,
  lem:finite-cap-chain-frontier-dichotomy,prop:overlapping-epsilon-necks,
  lem:adjacent-canonical-caps-form-component}
\notready
Let $X$ be connected and suppose every point of $X$ is a neck center or lies
in a $(C,\epsilon)$-cap core.  If $X$ meets a cap core, then there is a cap
$C_0$ meeting $X$ such that no finite balanced neck extension from $C_0$ has a
frontier point of $X$ in the core of a strictly larger cap.  The choice is
made along the ordered cap/neck incidence relation; it is not an arbitrary
maximal element of an unproved partial order.
\end{lemma}

\begin{proof}
Start with any cap whose core meets $X$ and follow only extensions through a
frontier point of $X$.  The overlap proposition gives a unique axial
orientation for each extension, so the extensions form a nested chain.  An
infinite sequence of strict enlargements would be the nested-cap chain
excluded by \ref{lem:no-infinite-nested-canonical-caps}; hence the extension
process stops after finitely many steps.  At its terminal cap, the finite-cap
frontier dichotomy says that every remaining frontier point is either a neck
center (which would extend the chain) or a distinct second-cap point $C'$.
If $C_0\subseteq C'$, then $C'$ is a strictly larger cap reachable by the same
frontier extension, contradicting terminality.  Hence $C_0\nsubseteq C'$ and
\ref{lem:adjacent-canonical-caps-form-component} applies.  This is exactly the
stated maximal selection and is the cap-selection step in Appendix~A.21.
\end{proof}

\begin{theorem}[Connected sets covered by canonical cap cores and neck centers]
\label{thm:cap-and-neck-cover-classification}
\dcref{MT:app:A.21}
\source{morgan-tian}
\notready
For every $C<\infty$ and all sufficiently small $\epsilon>0$, let $X$ be
a connected subset of a smooth Riemannian three-manifold $M$ without
boundary.  Suppose every
point of $X$ is either the center of an $\epsilon$-neck or lies in the core
of a $(C,\epsilon)$-cap.  Then at least one of the following holds:
\begin{enumerate}
\item $X$ lies in a component which is the union of two
      $(C,\epsilon)$-caps and is diffeomorphic to
      $S^3$, $\mathbb{RP}^3$, or
      $\mathbb{RP}^3\#\mathbb{RP}^3$;
\item $X$ lies in a doubly $C$-capped $\epsilon$-tube component of one
      of those three diffeomorphism types;
\item $X$ lies in one $(C,\epsilon)$-cap;
\item $X$ lies in a $C$-capped $\epsilon$-tube;
\item $X$ lies in an $\epsilon$-tube;
\item $X$ lies in an $\epsilon$-fibration component covered by
      $\epsilon$-necks.
\end{enumerate}
\end{theorem}

\begin{proof}
\uses{def:canonical-neck-cap-regions,lem:neck-chain-product,lem:balanced-neck-chain-extension,lem:infinite-neck-chain-frontier,lem:finite-cap-chain-frontier-dichotomy,thm:separating-neck-centers-form-tube,lem:nonseparating-neck-chain-or-bundle,lem:all-neck-component-tube-or-bundle,lem:neck-separation-type-locally-constant,lem:maximal-cap-selection,lem:adjacent-canonical-caps-form-component,prop:overlapping-epsilon-necks}
We first treat the case in which some $x_0\in X$ lies in the core of a cap.
Choose a cap $C_0$ containing $x_0$ using the ordered maximal-cap contract
\ref{lem:maximal-cap-selection}; this choice is not an appeal to an
unproved global maximality principle.  We now follow the component
of $X$ starting at $C_0$, keeping the cap orientation fixed so that the
negative end points toward its core.

There are three mutually exclusive possibilities.  If $X$ is disjoint from
$\operatorname{Fr}_M(C_0)$, then the connected set $X$, which meets $C_0$,
is contained in $C_0$; this is alternative (3).  Otherwise suppose every
point of $X\cap\operatorname{Fr}_M(C_0)$ is a neck center.  At a frontier
point of the positive end, the overlap proposition and the balanced-chain
extension lemma adjoin a neck (after reversing its axial coordinate if
necessary).  Iterating produces a finite or infinite balanced chain
  $N_1,N_2,\ldots$ attached to $C_0$.  A finite chain which stops contains
  $X$, by \ref{lem:finite-cap-chain-frontier-dichotomy}: a frontier neck would
  extend the chain, while a frontier cap is the second-cap case below.  If one end
is infinite, \ref{lem:infinite-neck-chain-frontier} says that this end has no
frontier; if the finite end cannot be extended, connectedness again puts all
of $X$ in the capped chain.  If both ends are infinite, the same lemma says
that the chain is an ambient component.  In every subcase the product lemma
identifies the neck union with $S^2\times(0,1)$, so we obtain alternative
(4).

It remains that a point of $X\cap\operatorname{Fr}_M(C_0)$ lies in the core
of another cap.  During the preceding extension, take the finite terminal
chain whose frontier contains that point (an infinite end cannot contribute a
frontier point).  Call the resulting capped tube $\widetilde C$ and the
second cap $C'$.  Maximality of $C_0$ gives $C_0\nsubseteq C'$, so
\ref{lem:adjacent-canonical-caps-form-component} shows that
$\widetilde C\cup C'$ is an entire component of $M$ containing $X$.
The intervening necks have product structure by
\ref{lem:neck-chain-product}; consequently this component is obtained by
gluing the compact core of $C_0$ to the compact core of $C'$ along one
two-sphere.  Each compact core is either a $3$-ball or
$\mathbb{RP}^3\setminus\operatorname{int}(B^3)$.  The four possible
gluings therefore reduce to $S^3$, $\mathbb{RP}^3$, and
$\mathbb{RP}^3\#\mathbb{RP}^3$.  Keeping the intervening product as a tube
gives alternative (2), while absorbing it into the two cap descriptions
gives alternative (1).

Finally assume that every point of $X$ is a neck center.  By
\ref{lem:neck-separation-type-locally-constant}, the central spheres are
either all separating or all nonseparating.  In the separating case,
\ref{thm:separating-neck-centers-form-tube} supplies a balanced chain whose
union contains $X$, giving alternative (5).  In the nonseparating case,
  \ref{lem:all-neck-component-tube-or-bundle} gives either the same product
  tube (again alternative (5)) or an $S^2$-bundle over $S^1$ (alternative
  (6)).  These cases exhaust the cap and neck hypotheses.
\end{proof}

\begin{theorem}[Topology of a manifold covered by canonical caps and necks]
\label{thm:canonical-neck-cap-topology}
\dcref{MT:app:A.25}
\source{morgan-tian}
\notready
Let $(M,g)$ be a connected smooth Riemannian three-manifold without boundary.
For fixed
$C<\infty$ and sufficiently small $\epsilon>0$, suppose each point of
$M$ is either the center of an $\epsilon$-neck or belongs to the core of a
$(C,\epsilon)$-cap.  At least one of the following source alternatives
occurs (their displayed diffeomorphism types are separated by compactness
and end structure):
\begin{enumerate}
\item $M$ is diffeomorphic to $S^3$, $\mathbb{RP}^3$, or
      $\mathbb{RP}^3\#\mathbb{RP}^3$, and is the union of two caps or a
      doubly capped $\epsilon$-tube;
\item $M$ is diffeomorphic to $\mathbb R^3$ or to
      $\mathbb{RP}^3\setminus\{p\}$, and is a cap or a capped
      $\epsilon$-tube;
\item $M$ is diffeomorphic to $S^2\times\mathbb R$ and is an
      $\epsilon$-tube;
\item $M$ is diffeomorphic to an $S^2$-bundle over $S^1$ and is an
      $\epsilon$-fibration.
\end{enumerate}
\end{theorem}

\begin{proof}
\uses{thm:cap-and-neck-cover-classification}
Apply \ref{thm:cap-and-neck-cover-classification} with $X=M$.  Since $M$
is a whole connected component, alternatives (1) and (2) give the first
type.  Alternatives (3) and (4) give the second type, using the two possible
cap-core topologies.  Alternative (5), together with the product definition
of a tube, gives the third type, and alternative (6) gives the fourth.
\end{proof}
